% source: https://github.com/pfradin/luadraw/discussions/255#discussioncomment-18029257 (pfradin, block 1)
\documentclass[border=2 mm]{standalone}
\usepackage[svgnames]{xcolor}
\usepackage[3d]{luadraw}
\usepackage{fouriernc}
\begin{document}
%%https://tex.stackexchange.com/questions/739936/fill-area-between-hyperbola-and-line-with-pgfplots
\begin{luadraw}{name=fill}
local ld = luadraw
local cpx = ld.cpx
local sqrt = math.sqrt
local Z, i = cpx.Z, cpx.I
local xmin, xmax, ymin, ymax = -18, 18, -2, 12
local g = ld.graph:new{window={xmin, xmax, ymin, ymax}, size={12,12}}  
g:Daxes({0,2,2}, {grid=true, gridstyle = "dashed",arrows="-Stealth", legend={"$x$","$y$"}})
g:Linewidth(12)
g:Linecap("round")
local f = function(x) return sqrt(3)/4*sqrt(256-x^2)  end
local h = function(x) return sqrt(100-x^2) end
g:Dinequalities({f,'>', h,'<'}, {view={-10,10,0,10}, draw_options="draw=none,fill=green,fill opacity=0.6"})
g:Dcartesian(f, {draw_options="blue"});
g:Dcartesian(h, {draw_options="violet"});
g:Show()
\end{luadraw}
\end{document} 
